% ============================================================
% 迭代三 Agent 优化工作记录（LaTeX 版）
% 编译方式：xelatex -> xelatex
% ============================================================
\documentclass[12pt,a4paper,onecolumn]{ctexart}

\usepackage[top=2.5cm,bottom=2.5cm,left=2.8cm,right=2.8cm]{geometry}
\usepackage{amsmath,amssymb}
\usepackage[table,xcdraw]{xcolor}
\usepackage{hyperref}
\usepackage{enumitem}
\usepackage{booktabs}
\usepackage{longtable}
\usepackage{array}
\usepackage{caption}
\usepackage{float}
\usepackage{tikz}
\usetikzlibrary{shapes.geometric, arrows.meta, positioning, fit, calc, backgrounds}
\emergencystretch=2em

\hypersetup{colorlinks=true, linkcolor=blue!60!black, urlcolor=blue!60!black}
\newcommand{\resizetikz}[2][\textwidth]{\resizebox{#1}{!}{#2}}

\title{\textbf{迭代三 Agent 优化工作记录}}
\author{软件工程课程项目组}
\date{2026年6月}

\begin{document}
\maketitle
\thispagestyle{empty}
\newpage
\tableofcontents
\newpage

%============================================================
\section{优化背景}
%============================================================

迭代二系统以一次性自动生成为主，用户输入主题后等待 Agent 自动搜索、记录笔记并生成综述。该方式自动化程度高，但存在过程黑盒、难以纠偏、中间产物不可编辑、评估和答辩可证明性不足等问题。

迭代三围绕“连续对话 + 交互式调研 + 自动模式 + 可观测性”展开，将系统升级为可交互、可追踪、可修订、可评估的学术调研助手。

%============================================================
\section{优化路线}
%============================================================

\begin{figure}[H]
\centering
\resizetikz{\begin{tikzpicture}[
  node distance=1.0cm,
  wave/.style={rectangle, draw=blue!70, fill=blue!6, rounded corners=2mm, text width=3.8cm, align=center, minimum height=0.9cm},
  arr/.style={-{Stealth}, thick}
]
  \node[wave] (w1) {第一波\\Session 与前端控制台};
  \node[wave, right=of w1] (w2) {第二波\\笔记编辑与修订};
  \node[wave, right=of w2] (w3) {第三波\\RAG 与 Paper Register};
  \node[wave, below=1.1cm of w3] (w4) {第四波\\自动模式、工具、Skills};
  \node[wave, left=of w4] (w5) {第五波\\分析阶段与写作注入};
  \node[wave, left=of w5] (w6) {第六波\\Skill Trace 与 Copilot 范围};

  \draw[arr] (w1) -- (w2);
  \draw[arr] (w2) -- (w3);
  \draw[arr] (w3) -- (w4);
  \draw[arr] (w4) -- (w5);
  \draw[arr] (w5) -- (w6);
\end{tikzpicture}}
\caption{迭代三优化路线}
\end{figure}

\begin{longtable}{p{2.5cm}p{5.5cm}p{5.5cm}}
\toprule
\textbf{阶段} & \textbf{主要问题} & \textbf{优化结果} \\
\midrule
Session 化 & 一次性运行，无法恢复 & 引入 Session 状态机与目录持久化。\\
交互式流程 & 用户无法干预关键词、论文和笔记 & 控制台提供分阶段审核与编辑。\\
RAG 笔记 & 笔记结构不稳定，PDF 利用不足 & 基于全文检索生成结构化笔记。\\
自动模式 & 分步操作对演示和快速使用成本较高 & 一键执行完整 Pipeline。\\
分析阶段 & 分析只是页面结果，未参与写作 & 分析结果注入 Writer。\\
可观测性 & Skill 是否生效难以证明 & 三阶段记录 \texttt{SKILL\_STATUS}。\\
知识共享 & 跨项目问答范围不可控 & Copilot 支持 \texttt{session\_ids}。\\
\bottomrule
\end{longtable}

%============================================================
\section{最终 Pipeline}
%============================================================

\begin{figure}[H]
\centering
\resizetikz{\begin{tikzpicture}[
  node distance=1.15cm,
  phase/.style={rectangle, draw=teal!70, fill=teal!6, rounded corners=2mm, text width=2.6cm, align=center, minimum height=0.85cm},
  artifact/.style={rectangle, draw=orange!70, fill=orange!7, rounded corners=1mm, text width=2.6cm, align=center, minimum height=0.75cm},
  arr/.style={-{Stealth}, thick},
  data/.style={-{Stealth}, dashed, gray, thick}
]
  \node[phase] (plan) {规划};
  \node[phase, right=of plan] (search) {搜索};
  \node[phase, right=of search] (notes) {笔记};
  \node[phase, right=of notes] (analysis) {分析};
  \node[phase, right=of analysis] (write) {综述};

  \node[artifact, below=1.0cm of plan] (a1) {关键词};
  \node[artifact, below=1.0cm of search] (a2) {论文列表};
  \node[artifact, below=1.0cm of notes] (a3) {笔记 Markdown};
  \node[artifact, below=1.0cm of analysis] (a4) {compare\\lineage\\gaps};
  \node[artifact, below=1.0cm of write] (a5) {综述草稿};

  \draw[arr] (plan) -- (search);
  \draw[arr] (search) -- (notes);
  \draw[arr] (notes) -- (analysis);
  \draw[arr] (analysis) -- (write);
  \draw[data] (plan) -- (a1);
  \draw[data] (search) -- (a2);
  \draw[data] (notes) -- (a3);
  \draw[data] (analysis) -- (a4);
  \draw[data] (write) -- (a5);
  \draw[data] (a4.east) .. controls +(0.8,-1.0) and +(-0.8,-1.0) .. node[below]{写作上下文} (write.south);
\end{tikzpicture}}
\caption{当前正式 Pipeline}
\end{figure}

%============================================================
\section{重点优化说明}
%============================================================

\subsection{Session 与状态机}
每个研究任务映射为一个 Session，包含 metadata、plan、papers、notes、analysis、draft、traces、chats 等目录。状态机覆盖从 planning 到 complete 的完整生命周期，并支持卡住状态修复。

\subsection{Web 解耦}
原先 Web 入口职责过重，后期拆分为 pages、session、agent、chat、conversation、draft、admin 等路由模块。通过 \texttt{deps.py} 延迟注入共享依赖，降低循环导入风险。

\subsection{RAG 笔记生成}
笔记阶段使用全文切块与向量检索，针对背景、方法、实验、结果、分析、局限等不同部分提取相关上下文。Embedding 失败时回退 BM25，PDF 失败时回退摘要。

\subsection{分析阶段并入主流程}
分析阶段从独立页面能力升级为正式 Pipeline 节点。自动模式中，系统先生成分析再写综述。Writer 读取 \texttt{analysis\_results.json}，将用户编辑后的分析文档作为额外上下文。

\subsection{Skills 可证明}
Skills 支持 search、notes、write 三阶段定制。系统在每个阶段都写入 \texttt{SKILL\_STATUS} trace，记录 skill id、标题、是否加载、是否 fallback 和原因。

\subsection{全局 Copilot 范围选择}
全局 Copilot 原先默认扫描所有 Session。当前支持前端多选 Session，后端知识问答接口根据 \texttt{session\_ids} 限定检索范围。

%============================================================
\section{能力雷达图语义}
%============================================================

\begin{figure}[H]
\centering
\resizetikz[0.72\textwidth]{\begin{tikzpicture}[
  scale=1.0,
  axis/.style={gray!70, thin},
  label/.style={font=\small},
  poly/.style={fill=blue!25, draw=blue!70, thick, opacity=0.75}
]
  \def\r{3}
  \foreach \i/\name in {
    90/交互性,
    30/自动化,
    -30/可观测,
    -90/可评估,
    -150/可维护,
    150/可扩展
  } {
    \draw[axis] (0,0) -- (\i:\r);
    \node[label] at (\i:\r+0.45) {\name};
  }
  \foreach \rad in {1,2,3} {
    \draw[axis] (90:\rad) -- (30:\rad) -- (-30:\rad) -- (-90:\rad) -- (-150:\rad) -- (150:\rad) -- cycle;
  }
  \draw[poly] (90:2.8) -- (30:2.7) -- (-30:2.9) -- (-90:2.4) -- (-150:2.6) -- (150:2.5) -- cycle;
  \node at (0,-3.7) {迭代三后系统能力分布示意};
\end{tikzpicture}}
\caption{迭代三优化后的能力覆盖}
\end{figure}

%============================================================
\section{任务点对齐}
%============================================================

\begin{longtable}{p{3.4cm}p{7.0cm}p{2.6cm}}
\toprule
\textbf{任务点} & \textbf{当前实现} & \textbf{判断} \\
\midrule
连续对话 & Session 内多会话聊天、历史保存、上下文压缩 & 已实现 \\
交互式调研 & 关键词、论文、笔记、分析、综述均可人工介入 & 已实现 \\
自动执行 & 自动模式一键执行完整流程 & 已实现 \\
混合论文来源 & Agent 搜索、链接添加、PDF 上传 & 已实现 \\
中间产物编辑 & 笔记、分析卡片、综述均可 Markdown 编辑 & 已实现 \\
运行轨迹 & ReAct trace、自动流程 trace、Skill trace & 已实现 \\
深度分析 & compare/lineage/gaps 并注入写作 & 已实现 \\
跨 Session 知识 & 全局 Copilot + 可选 Session 范围 & 已实现 \\
评估支撑 & evaluation 目录、数据、脚本与 fallback & 基本具备 \\
\bottomrule
\end{longtable}

%============================================================
\section{风险与后续建议}
%============================================================

\begin{enumerate}
  \item 增强 trace 面板摘要，展示工具调用次数、错误类型统计、论文收录数。
  \item 在分步模式论文卡片中展示是否有笔记、是否被综述引用、是否适合用于对比/脉络/空白。
  \item 补全 evaluation 数据集，增加自动模式、分析注入、Skill trace、Copilot 范围选择的专项样例。
  \item 为前端核心流程增加轻量端到端回归测试。
\end{enumerate}

\end{document}
